\section{Conclusione}
\label{sec:conclusione}

\subsection{Riepilogo dei Risultati}

Il progetto \textbf{Videogame Semantic Search} ha prodotto un sistema completo
di ricerca semantica che integra:

\begin{itemize}[nosep]
  \item \textbf{Un'ontologia OWL 2} espressiva e ben strutturata, evoluta da 11 a
    78 classi (61 primitive + 17 definite), con 31 object property, 26 datatype
    property, 13 coppie inverse, 3 property chain, 2 propriet\`a transitive,
    12 blocchi di disjointness e allineamento esterno a Wikidata e Dublin Core;
  \item \textbf{Un sistema di reasoning a due livelli} (OWL 2 RL + 21 regole
    CONSTRUCT) funzionante end-to-end, con 0 falsi \texttt{owl:sameAs} dopo
    il fix del bug critico degli assiomi DL-only;
  \item \textbf{Un knowledge graph popolato} da Wikidata con oltre 1.1 milioni
    di triple (dataset completo 2010--2026) e $\sim$745k triple nel sottoinsieme
    demo (2020--2026, $\sim$68.700 giochi);
  \item \textbf{Un'interfaccia web interattiva} con ricerca in linguaggio naturale,
    agente SPARQL intelligente, grafo 2D navigabile e filtri per tipo;
  \item \textbf{Performance ottimizzate} grazie al passaggio da rdflib a pyoxigraph
    (5--62$\times$ speedup sulle query, caricamento in $<$3s, RAM $\sim$250 MB).
\end{itemize}

\subsection{Possibili Estensioni Future}

Il sistema offre numerose direzioni di sviluppo futuro:

\begin{enumerate}[nosep]
  \item \textbf{Reasoner DL completo} --- integrare HermiT o Pellet tramite
    \texttt{owlready2} per ragionare sui costrutti DL-only (datatype restriction,
    complementOf, maxCardinality) senza bisogno delle regole CONSTRUCT
    compensate. Questo abiliterebbe anche la rilevazione automatica di
    inconsistenze tra le istanze;
  \item \textbf{Validazione delle istanze con SHACL} --- definire SHACL shapes per
    validare la conformit\`a delle istanze (es. ogni VideoGame deve avere
    almeno un \texttt{gameName}, le date devono essere nel formato corretto);
  \item \textbf{SWRL rules reali} --- sostituire le 21 regole CONSTRUCT con
    autentiche SWRL rules, sfruttando un reasoner che le supporti nativamente;
  \item \textbf{Estensione della pipeline Wikidata} --- aggiungere query per
    recuperare \texttt{hasSteamAppId}, \texttt{budgetTier}, \texttt{playerCount}
    e la relazione \texttt{sequelOf}, attualmente non popolate dai dati Wikidata;
  \item \textbf{Internazionalizzazione del frontend} --- supporto multilingua
    per l'interfaccia e le query in linguaggio naturale (inglese, francese, \dots{});
  \item \textbf{Visualizzazione avanzata delle classi definite} --- nel grafo
    interattivo, evidenziare con badge o colori speciali i nodi che sono stati
    classificati automaticamente dal reasoner (es. bordino dorato per
    \texttt{HighlyRatedGame}, icona controller per \texttt{MultiplayerGame});
  \item \textbf{SPARQL endpoint pubblico} --- esporre un endpoint SPARQL
    standard (RESTful o tramite Triplestore) per consentire query federate
    da applicazioni esterne.
\end{enumerate}

Il progetto dimostra come le tecnologie del Semantic Web --- ontologie formali,
reasoning automatico, knowledge graph e query SPARQL --- possano essere integrate
in un'applicazione web moderna e performante, offrendo un'esperienza di ricerca
semantica che va ben oltre la semplice ricerca per parole chiave.
